% Source: round12/stability/prime_rates.tex
\section{Quantitative stability counterexamples from ordinary prime gaps}
\label{s26:sec:stability-prime-gaps}
This section replaces the irrational-rotation subsequence with a construction
at every sufficiently large dyadic seed order. No prime progression is
needed. We use the established Baker--Harman--Pintz theorem that
$[x-x^{21/40},x]$ contains a prime for all sufficiently large $x$:
\emph{The difference between consecutive primes, II} (2001), Theorem 1,
\url{https://www.cs.umd.edu/~gasarch/BLOGPAPERS/BakerHarmanPintz.pdf}.
We do not claim $21/40$ is the best current prime-gap exponent.
The separate conditional improvement uses ordinary RH, not GRH for
arithmetic progressions. No new Hadamard orders are asserted.

\begin{lemma}[No residue-class restriction is needed]
\label{s26:lem:stability-any-prime}
For every odd prime $p$, a Hadamard matrix of order $k=2(p+1)$ exists.
\end{lemma}
\begin{proof}
For $p\equiv1\pmod4$, this is classical Paley type II. For
$p\equiv3\pmod4$, Paley type I gives order $p+1$, which doubles to
$2(p+1)$. Both constructions originate in Paley,
\emph{On Orthogonal Matrices} (1933),
\url{https://doi.org/10.1002/sapm1933121311}.
\end{proof}

\paragraph{Formalization.}
The lemma is formalized from scratch in the companion Lean project.
For type II, a symmetric conference matrix $C$ of order $p+1$
satisfies $C^2=pI$, with zero diagonal and sign off-diagonal entries.
The block matrix
\[
 \begin{pmatrix} C+I&C-I\\ C-I&-C-I\end{pmatrix}
\]
has sign entries and Gram matrix $2(p+1)I$.
The formal theorem uses mathlib's standard Hadamard predicate and is
specified independently without new definition holes. This formalizes
a classical input, not a new existence theorem. The ordinary-prime-gap transfer below is formalized separately.
The unconditional BHP and RH applications retain their own proof obligations.


\begin{theorem}[An ordinary-prime-gap transfer]
\label{s26:thm:stability-prime-transfer}
Suppose $0<\theta<1$ and every sufficiently large interval
$[x-x^\theta,x]$ contains a prime. For every sufficiently large power
of two $h$, there is an explicitly searchable known Hadamard order
$k<h$ with
\[
 h-k=\Theta(h^\theta),\qquad k/h\to1.
\]
At the known Hadamard order $n=kh$, there is an exact unbiased
pairwise-independent uniform $h^2$-atom law $\mu_h$ on $n-1$ bits,
a strict local collision maximum, for which
\begin{align*}
 C(\mu_h)&=1/h^2, & C_{\rm global}&=1/n,\\
 C(\mu_h)/C_{\rm global}&=1-\Theta(n^{-(1-\theta)/2}),&
 \operatorname{dist}_{\rm TV}(\mu_h,\mathcal H_n)
 &=1-1/k=1-\Theta(n^{-1/2}),\\
 1/n-C(\mu_h)&=\Theta(n^{-(3-\theta)/2}).
\end{align*}
All moments are exact; there is no approximate-isotropy assumption.
\end{theorem}
\begin{proof}
Let $X=h/2$, $A=\lceil X^\theta\rceil$, and $y=X-A$.
Choose the largest prime $p\le y$. The input theorem, applied at $y$,
gives $p\ge y-y^\theta$. Set $k=2(p+1)$, so
\[
 2A-2\le g:=h-k\le2A+2y^\theta-2\le4A-2.
\]
Thus $g=\Theta(h^\theta)$, and eventually
$\log_2h\le g<h/2$. The lemma supplies $H_k$, hence $H_n$ by tensoring
with Walsh order $h$.

Choose an omitted-slope set $R\subset\mathbb F_h\setminus\{0\}$ of
size $g$ containing an additive basis, and put
$P=\mathbb F_h\setminus R$. For uniform $A_0,B_0\in\mathbb F_h$, observe
all nonzero additive characters of each seed block
$x_t=A_0+tB_0$, $t\in P$, and the mixed bits $\chi_e(x_0+x_t)$,
$t\in P\setminus\{0\}$. This is the spanning-omitted-slopes construction
of Theorem~\ref{s25:thm:design-spanning-slopes}: it has $h^2$ distinct atoms,
exact moments, and is a strict local maximum because two block character
spaces generate every seed character through products. Its orthogonality
graph is $h$ disjoint $h$-cliques; their order-$k$ block completions give
exact distance $1-1/k$ from all Hadamard mixtures. Collision is $1/h^2$,
whereas the known full optimum is $1/(kh)$.
Finally $n\asymp h^2$, and
\[
 1-\frac{C(\mu_h)}{1/n}=\frac gh,\qquad
 \frac1n-C(\mu_h)=\frac{g}{kh^2}.
\]
These identities give all rates.
\end{proof}

\paragraph{Formalization of the ordinary-prime-gap transfer.}
The independent \texttt{PrimeChallenge} verifies the complete conditional
transfer from the stated ordinary prime intervals, using actual integral
Paley matrices and one family of exact finite laws. It proves the relative
and absolute collision deficit rates and the inverse-order distance rate
with real powers of the full order $n=kh$. The same laws have the exact
support and collision identities, strict local optimality, and an attained
sharp distance from all Hadamard mixtures. Only the prime-gap hypothesis
is assumed. The BHP and RH inputs are separate analytic obligations.

\begin{corollary}[Unconditional polynomial rates]
\label{s26:cor:stability-bhp-rate}
Unconditionally the theorem holds with $\theta=21/40$. In particular
\[
 1-C(\mu_h)/(1/n)=\Theta(n^{-19/80}),\qquad
 1/n-C(\mu_h)=\Theta(n^{-99/80}),\qquad
 \operatorname{dist}_{\rm TV}(\mu_h,\mathcal H_n)=1-\Theta(n^{-1/2}).
\]
The order selection is a finite descending prime search below
$h/2-\lceil(h/2)^{21/40}\rceil$; the asymptotic theorem does not claim
a numerical value for the Baker--Harman--Pintz threshold.
\end{corollary}

\subsection{A separate improvement under ordinary RH}
Dudek proves on RH that for every $x\ge2$ there is a prime in
$(x-\frac4\pi\sqrt{x}\log x,x]$:
\emph{On the Riemann hypothesis and the difference between primes},
\url{https://arxiv.org/abs/1402.6417}. This is a primary, explicit
ordinary-prime result; no arithmetic-progression hypothesis is used.
\begin{corollary}[Conditional RH rates]\label{s26:cor:stability-rh-rate}
Assuming RH, the same construction can be chosen at every sufficiently
large dyadic $h$ with
\[
 h-k=\Theta(\sqrt h\log h),\quad
 1-C(\mu_h)/(1/n)=\Theta(n^{-1/4}\log n),\quad
 1/n-C(\mu_h)=\Theta(n^{-5/4}\log n),
\]
while the exact distance remains $1-\Theta(n^{-1/2})$.
\end{corollary}
\begin{proof}
Use $A=\lceil\frac4\pi\sqrt X\log X\rceil$ and $y=X-A$ in the
preceding argument. The function $\frac4\pi\sqrt x\log x$ is increasing
for $x\ge2$, so the prime supplied at $y$ gives
$2A-2\le h-k\le4A-2$. The remaining proof is unchanged.
\end{proof}

\paragraph{Entropy and the scale of a possible stability theorem.}
Every Renyi entropy, including Shannon and min-entropy, of the uniform
source equals $\log h^2$. The minimum among feasible laws at known order
$n$ is $\log n$, by the atom bound and the Hadamard law. Thus all these
entropy excesses are simultaneously
$\log(h/k)=\Theta(n^{-19/80})$ in the unconditional construction,
and $\Theta(n^{-1/4}\log n)$ under RH, although distance tends to one.
These examples refute a dimension-free stability modulus based on relative
collision deficit or entropy excess. They also quantify absolute-deficit
counterexamples, at the exponents above. They do not exclude much smaller
dimension-dependent absolute collision windows, such as the earlier
order-$n^{-4}$ extraction condition. All positive existence inputs here
are classical known Hadamard constructions.

\begin{corollary}[Only sublinearly many extra atoms]
\label{s26:cor:stability-support-excess}
The unconditional source support size satisfies
\[
 |\operatorname{supp}\mu_h|=n+\Theta(n^{61/80}).
\]
Under RH it can instead be $n+\Theta(n^{3/4}\log n)$, while its exact
distance from all Hadamard mixtures still tends to one.
\end{corollary}
\begin{proof}
The atom count is $h^2$ and $h^2-n=h(h-k)$. Substitute the corresponding
prime-gap scale and $n\asymp h^2$.
\end{proof}


% Formalization scope: E26.03
\paragraph{Lean status (\texttt{E26.03}).} The complete principal statement has not been formalized. Supporting lemmas or computational certificates do not establish the whole claim.

% Formalization scope: E26.04
\paragraph{Lean status (\texttt{E26.04}).} The complete principal statement has not been formalized. Supporting lemmas or computational certificates do not establish the whole claim.

% Formalization scope: E26.05
\paragraph{Lean status (\texttt{E26.05}).} The complete principal statement has not been formalized. Supporting lemmas or computational certificates do not establish the whole claim.
